% page 474 du fac-simile (image p-473 du scan)
% bloc 157.5 x 233.3 mm ; marge gauche 8.0 mm
\pgc{-12.700mm}{0.000mm}{
\l{\cel{3}466}
\l{}
\l{\sou{propulsar}. (\sou{trans,}\cel{1}ad.)(\sou{mekaniko}). (aludante helico, palet-}
\l{\cel{3}roto, e c.). Pulsar adavan, avance. - DEFIS.}
\l{}
\l{prorogar. I. (\sou{yurocienco}) Ajornar oficale a dato plu o min}
\l{\cel{3}fora. - II. (\sou{aludante asemblitaro}). Rezolvar cesar quik}
\l{\cel{3}la +seanco, la +sesiono, til epoko determinato. - DEFIS.}
\l{}
\l{\sou{prosektoro}. (\sou{anat.}) Helpisto qua preparas, por la profesoro}
\l{\cel{3}di medicino, la disseko-peci. -}
\l{}
\l{\sou{prosenkimo.}\cel{2}(\sou{bot.)}\cel{2}Tisuo ek flori. -}
\l{}
\l{\sou{prosilogismo}.\cel{2}(\sou{logiko}). Katenatra dispozeso di la silogismo,}
\l{\cel{3}en qua la konkluzo di la unesma divenas la premiso di la}
\l{\cel{3}duesma - e tale pluse. -}
\l{}
\l{\sou{proskriptar}. (\sou{trans.}). (En\cel{1}\sou{Roma, antique}). Imperar la ocido}
\l{\cel{3}(di ulu) sen formi judiciala, per nur enskribo sur afisho-}
\l{\cel{3}texto publikigita. - II. Kondamnar ad ocideso granda-}
\l{\cel{3}quante. - DEFIRS.}
\l{}
\l{\sou{prospekto}.\cel{2}Anuncilo, generale imprimura, pri verko qua bal-}
\l{\cel{3}de publikigesos, o pri institucuro por publiko, o pri varo,}
\l{\cel{3}e c., quan onu disdonas por atraktar la kompronti, la}
\l{\cel{3}klienti, eventuala.DEFIRS.}
\l{}
\l{\sou{prosperar.}\cel{2}(\sou{netrans.}) Favoresar da fortuno, sucesar, standar}
\l{\cel{3}plena-vigoro. - DEFIS.}
\l{}
\l{\sou{prostato}. (\sou{anat.)}\cel{3}Glando situita (che la viro) ye la junto-}
\l{\cel{3}punto di la veziko-kolo kun la uretro. - DEFIS.}
\l{}
\l{\sou{prostatito}. (\sou{patol.)}\cel{2}Influro di la prostato. - DEFIS.}
\l{}
\l{\sou{prosternar}. (\sou{netrans.})\cel{2}Kushar su, kun sua ventro vers sulo,}
\l{\cel{3}kun sua vizajo kontre sulo (od inklinar su til kande sua}
\l{\cel{3}fronto tushas sulo) koram ulu, kom signo di respekto.-}
\l{\cel{3}EFIS.}
\l{}
\l{\sou{prostezo}. (\sou{gram.}) Adjunto di un litero, di un silabo, komence}
\l{\cel{3}di vorto. (kom ex.:\cel{1}\sou{ica}\cel{1}vice ca, en L.\cel{1}\sou{iscala}\cel{1}vice\cel{1}\sou{scala}\cel{1};}
\l{\cel{3}en F.\cel{1}\sou{lierre}\cel{1}vice\cel{1}\sou{ierre}). - DEFIS.}
\l{}
\l{\sou{prostitucar}. (\sou{trans.}) I. (\sou{aludante muliero}). Koitar kun irgu}
\l{\cel{3}po pekunio. II. (\sou{metaf.)}\cel{2}\sou{Prostitucar la judicio:}\cel{1}sakrifi-}
\l{\cel{3}kar lu pro la expekto di pekunio, di favori, e c.\cel{1}\sou{Prosti}-}
\l{\cel{3}\sou{tucar sua plumo}\cel{1}: vendar sua skribo, sua redakto, po pe-}
\l{\cel{3}kunio, ofico, e c.\cel{2}\sou{Prostitucar amikeso}\cel{1}: igar irgu sua}
\l{\cel{3}amiko, expekte di profito. - DEFIRS.}
\l{}
\l{prostracar. (\sou{netrans.}) Esar abatita extreme, mentale, spiri-}
\l{\cel{3}tale, o korpale. - DEFIRS.}
}
